\section{Materials and Methods}

\subsection{Dataset}

The KiTS23 (Kidney Tumor Segmentation Challenge 2023) dataset \cite{kits23} provides contrast-enhanced abdominal CT volumes in NIfTI format with expert annotations for three classes: background (label 0), kidney (label 1), and tumor/cyst (label 2). The merged tumor/cyst class follows the challenge convention, reflecting the clinical difficulty of distinguishing these pathologies on single-phase CT. NephroVision evaluates on a project-defined held-out subset of 64 independent cases, selected before model development and never used during training or validation. All reported metrics refer to this 64-case evaluation subset, not an official KiTS23 leaderboard submission.

\subsection{Preprocessing}

All CT volumes undergo identical preprocessing:
\begin{itemize}
    \item \textbf{HU clipping:} Intensities clipped to $[-200, 300]$ Hounsfield Units to focus on soft-tissue contrast while suppressing bone, air, and extreme outliers.
    \item \textbf{Normalization:} Min--max scaling to $[0, 1]$: $x_{\text{norm}} = (x + 200) / 500$.
\end{itemize}

This preprocessing standardizes input distributions across scanners and acquisition protocols. Identical preprocessing is applied during training and inference to ensure consistency.

\subsection{Model Architecture}

The segmentation model is a 3D U-Net \cite{unet3d}, selected for its proven effectiveness in volumetric biomedical segmentation. The encoder-decoder architecture with skip connections processes volumetric patches and produces per-voxel class predictions. Three-dimensional convolutions capture inter-slice spatial context essential for kidney morphology, avoiding limitations of 2D slice-by-slice approaches. The model outputs three classes: background, kidney, and tumor/cyst.

\subsection{Training and Validation}

The model was trained on a subset of KiTS23 cases with a held-out validation set for hyperparameter tuning and checkpoint selection. Training used standard deep learning frameworks with GPU acceleration. Checkpoint selection was based on validation Dice score, with the best-performing model retained for final evaluation. Detailed training hyperparameters, loss functions, and optimization settings are documented in the experiment log.

\subsection{Inference Strategy}

Full CT volumes exceed GPU memory capacity. Sliding-window inference processes volumes with $64 \times 192 \times 192$ voxel patches and $50\%$ overlap in all three dimensions. Overlapping predictions are combined using averaging in overlap regions, producing smooth transitions across patch boundaries.

Test-time augmentation (TTA) applies 8 flip combinations (all combinations of flips along $x$, $y$, $z$ axes). Each augmented version is processed independently, predictions are un-flipped and averaged voxel-wise. TTA reduces boundary noise and improves segmentation consistency at approximately $8\times$ inference cost, acceptable for final evaluation.

\subsection{Post-Processing}

Connected-component filtering removes spurious detections:
\begin{itemize}
    \item \textbf{Kidney:} Components $< 5000$ voxels removed.
    \item \textbf{Tumor/cyst:} Components $< 100$ voxels removed.
\end{itemize}

The conservative tumor threshold preserves small true lesions while removing pixel-level noise. This prioritizes detection sensitivity over precision, consistent with decision-support use: physician review identifies and dismisses false positives, while false negatives (missed tumors) are more consequential.

\subsection{Output Generation}

The pipeline produces:
\begin{itemize}
    \item \textbf{Segmentation mask:} 3D volume with per-voxel class labels (background, kidney, tumor/cyst).
    \item \textbf{Volumetric statistics:} Kidney volume (ml), tumor/cyst volume (ml), tumor detection status.
    \item \textbf{3D visualization:} Browser-based interactive mesh viewer with rotate/zoom/pan and class toggling.
\end{itemize}

\begin{table}[!t]
\caption{System Configuration}
\label{tab:pipeline_config}
\centering
\begin{tabular}{@{}ll@{}}
\toprule
\textbf{Property} & \textbf{Value} \\
\midrule
Dataset & KiTS23 (2023) \\
Input format & NIfTI CT volumes \\
Classes & Background, Kidney, Tumor/Cyst \\
Test set & 64 held-out cases \\
\midrule
HU clipping & $[-200, 300]$ \\
Normalization & $[0, 1]$ (min--max) \\
\midrule
Model & 3D U-Net \\
Inference window & $64 \times 192 \times 192$ \\
Overlap & $50\%$ \\
TTA & 8 flip combinations \\
\midrule
Kidney blob threshold & $< 5000$ voxels \\
Tumor blob threshold & $< 100$ voxels \\
\bottomrule
\end{tabular}
\end{table}


Table~\ref{tab:pipeline_config} summarizes the system configuration.

\begin{figure}[!t]
\centering
\fbox{\parbox{0.8\linewidth}{\centering\vspace{1.5cm}\textbf{[Placeholder]}\\System Pipeline Diagram\\\small NIfTI Input $\rightarrow$ Preprocessing $\rightarrow$ 3D U-Net $\rightarrow$ Sliding-Window Inference $\rightarrow$ TTA $\rightarrow$ Post-processing $\rightarrow$ 3D Visualization\vspace{1.5cm}}}
\caption{NephroVision pipeline. CT volumes are preprocessed, segmented by 3D U-Net with sliding-window inference and test-time augmentation, post-processed to remove false positives, and visualized through a browser-based interface.}
\label{fig:pipeline}
\end{figure}
